% lecture09.tex — Properties of Brownian Motion

\section{Almost Sure Nowhere Differentiability of Brownian Motion}

Recall the upper and lower right derivatives.

\begin{definition}{Upper and Lower Right Derivatives}{right-derivs}
    Let $f \colon [0,\infty) \to \R$. Define the \emph{upper right derivative} and \emph{lower right derivative} of $f$ at $t$ by
    \[
        D^* f(t) = \limsup_{h \downarrow 0} \frac{f(t+h) - f(t)}{h}, \qquad
        D_* f(t) = \liminf_{h \downarrow 0} \frac{f(t+h) - f(t)}{h}.
    \]
\end{definition}

\begin{theorem}{Nowhere Differentiability of Brownian Motion (Paley, Wiener, Zygmund, 1933)}{pwz-nondiff}
    Brownian motion is nowhere differentiable almost surely. More precisely, almost surely, for all $t \geq 0$, either $D^* B(t) = +\infty$ or $D_* B(t) = -\infty$ (or both).
\end{theorem}

\begin{proof}
    Suppose for the sake of contradiction that there exists $t_0 \in [0,1]$ such that
    \[
        -\infty < D_* B(t_0) \leq D^* B(t_0) < \infty.
    \]
    Then
    \[
        \limsup_{h \downarrow 0} \frac{\abs{B(t_0 + h) - B(t_0)}}{h} < \infty
    \]
    and $\{B_t : t \in [0,2]\}$ is bounded (by continuity of paths on a compact set).

    This implies there exists $M < \infty$ such that
    \[
        \sup_{h \in [0,1]} \frac{\abs{B(t_0 + h) - B(t_0)}}{h} \leq M.
    \]
    We will show this event has probability $0$.

    \medskip
    \textbf{Fix $M > 0$.} For $n \geq 1$, consider the dyadic partition $\{k/2^n : k = 0, 1, \ldots, 2^n\}$ of $[0,1]$.

    If $t_0 \in [k/2^n, (k+1)/2^n)$ for some $k$, then for $j = 1, 2, 3$:
    \[
        \abs{B\!\left(\frac{k+j}{2^n}\right) - B\!\left(\frac{k+j-1}{2^n}\right)}
        \leq \abs{B\!\left(\frac{k+j}{2^n}\right) - B(t_0)} + \abs{B\!\left(\frac{k+j-1}{2^n}\right) - B(t_0)}.
    \]

    Since $\abs{B(t_0 + h) - B(t_0)} \leq M h$ for $h \in [0,1]$, and $0 \leq \frac{k+j}{2^n} - t_0 \leq \frac{j+1}{2^n}$ and $0 \leq \frac{k+j-1}{2^n} - t_0 \leq \frac{j}{2^n}$, we get:
    \begin{itemize}
        \item For $j = 1$:
        \[
            \abs{B\!\left(\frac{k+1}{2^n}\right) - B\!\left(\frac{k}{2^n}\right)} \leq \frac{3M}{2^n}.
        \]
        \item For $j > 1$:
        \[
            \abs{B\!\left(\frac{k+j}{2^n}\right) - B\!\left(\frac{k+j-1}{2^n}\right)} \leq \frac{M(2j+1)}{2^n}.
        \]
    \end{itemize}

    Define the event
    \[
        \Omega_{n,k} = \left\{ \abs{B\!\left(\frac{k+j}{2^n}\right) - B\!\left(\frac{k+j-1}{2^n}\right)} \leq \frac{M(2j+1)}{2^n} \text{ for } j = 1, 2, 3 \right\}.
    \]

    \textbf{Bounding $\Prob(\Omega_{n,k})$:} By independence of the increments,
    \[
        \Prob(\Omega_{n,k}) \leq \prod_{j=1}^{3} \Prob\!\left[\abs{B\!\left(\frac{k+j}{2^n}\right) - B\!\left(\frac{k+j-1}{2^n}\right)} \leq \frac{M(2j+1)}{2^n}\right].
    \]
    Since $B\!\left(\frac{k+j}{2^n}\right) - B\!\left(\frac{k+j-1}{2^n}\right) \sim N(0, 1/2^n)$, we write this increment as $\frac{1}{2^{n/2}} Z$ where $Z \sim N(0,1)$. Then
    \[
        \Prob\!\left[\abs{\frac{1}{2^{n/2}} Z} \leq \frac{M(2j+1)}{2^n}\right]
        = \Prob\!\left[\abs{Z} \leq \frac{M(2j+1)}{2^{n/2}}\right]
        \leq \frac{M(2j+1)}{2^{n/2}} \cdot \sqrt{\frac{2}{\pi}} \leq \frac{M(2j+1)}{2^{n/2}}.
    \]
    Using $2j+1 \leq 7$ for $j = 1,2,3$:
    \[
        \Prob(\Omega_{n,k}) \leq \left(\frac{7M}{2^{n/2}}\right)^3.
    \]

    \textbf{Union bound:} Summing over $k = 0, 1, \ldots, 2^n - 1$:
    \[
        \Prob\!\left[\bigcup_{k=0}^{2^n - 1} \Omega_{n,k}\right]
        \leq 2^n \cdot \left(\frac{7M}{2^{n/2}}\right)^3
        = (7M)^3 \cdot 2^n \cdot 2^{-3n/2}
        = \frac{(7M)^3}{2^{n/2}}.
    \]
    This is summable in $n$.

    \medskip
    \textbf{Applying Borel--Cantelli:} Since $\sum_{n=1}^{\infty} \frac{(7M)^3}{2^{n/2}} < \infty$, by the first Borel--Cantelli lemma,
    \[
        \Prob\!\left[\exists\, t_0 \in [0,1] \text{ with } \sup_{h \in [0,1]} \frac{\abs{B(t_0+h) - B(t_0)}}{h} \leq M\right]
        \leq \Prob\!\left[\bigcup_{k} \Omega_{n,k} \text{ for infinitely many } n\right] = 0.
    \]

    Finally, let $M \uparrow \infty$ (through integers):
    \[
        \Prob\!\left[\exists\, t_0 \in [0,1] \text{ with } D^* B(t_0) < \infty \text{ and } D_* B(t_0) > -\infty\right]
        \leq \sum_{M=1}^{\infty} 0 = 0. \qedhere
    \]
\end{proof}

%────────────────────────────────────────────────────────────────
\section{Markov Property of Brownian Motion}
%────────────────────────────────────────────────────────────────

Recall: $\{X_n : n \geq 1\}$ is a Markov chain if
\[
    \Prob[X_{n+1} \in A \mid X_0, \ldots, X_n] = \Prob[X_{n+1} \in A \mid X_n].
\]

\textbf{Recasting:} Equivalently,
\[
    \Prob[X_{n+1} = x_{n+1},\, X_{n-k} = x_{n-k} \mid X_n = x_n]
    = \Prob[X_{n+1} = x_{n+1} \mid X_n = x_n] \cdot \Prob[X_{n-k} = x_{n-k} \mid X_n = x_n].
\]

\textbf{Moral:} Given the present, the future and the past are independent.

%────────────────────────────────────────────────────────────────
\section{Filtrations}
%────────────────────────────────────────────────────────────────

\begin{definition}{Filtration}{filtration}
    Let $(\Omega, \F, \Prob)$ be a probability space. A \emph{filtration} $\{\F_t : t \geq 0\}$ is a collection of $\sigma$-algebras such that
    \[
        \F_s \subseteq \F_t \quad \text{for all } s \leq t.
    \]
\end{definition}

\begin{example}{Natural Filtration of Brownian Motion}{natural-filtration}
    Define $\F_t^0 = \sigma(\{B_s : s \leq t\})$.
\end{example}

\begin{lemma}{Shifted Brownian Motion}{shifted-bm}
    Let $\{B_t : t \geq 0\} \sim \mathrm{sBM}$ and define $W_t = B_{t+s} - B_s$. Then $\{W_t : t \geq 0\}$ is a standard Brownian motion independent of $\{B_t : 0 \leq t \leq s\}$.

    Formally: for all $A \in \B(C([0,s]))$ measurable and $B \in \B(C([0,\infty)))$ measurable,
    \[
        \Prob\!\left[(W_t)_{t \geq 0} \in B,\; (B_t)_{t \leq s} \in A\right]
        = \Prob\!\left[(B_t)_{t \leq s} \in A\right] \cdot \Prob\!\left[(W_t)_{t \geq 0} \in B\right].
    \]
\end{lemma}

\begin{proof}
    Let $\{X_t : 0 \leq t \leq s\}$ and $\{Y_t : t \geq 0\}$ be independent Brownian motions. Define
    \[
        B_t^* =
        \begin{cases}
            X_t & \text{if } 0 \leq t \leq s, \\
            X_s + Y_{t-s} & \text{if } t > s.
        \end{cases}
    \]
    Then $\{B_t^* : t \geq 0\}$ is a standard Brownian motion (by the gluing construction).

    Note that $(B_t^*)_{t \leq s} = (X_t)_{t \leq s}$ and $B_{t+s}^* - B_s^* = Y_t$, so the shifted process equals $\{Y_t : t \geq 0\}$, which is independent of $(X_t)_{t \leq s}$ by construction.

    By independence of $X$ and $Y$, for measurable $A$ and $B$:
    \[
        \Prob\!\left[(Y_t)_{t \geq 0} \in B,\; (X_t)_{t \leq s} \in A\right]
        = \Prob\!\left[(X_t)_{t \leq s} \in A\right] \cdot \Prob\!\left[(Y_t)_{t \geq 0} \in B\right].
    \]

    Since $\{B_t^* : t \geq 0\} \stackrel{d}{=} \{B_t : t \geq 0\}$, the same factorization holds for $(B_t)$. It remains to note that $\sigma(\{f : f(t_i) \in A_i, i = 1, \ldots, k\}) = \B(C([0,s]))$ (by a $\pi$-$\lambda$ argument), completing the proof.
\end{proof}
